% Independent mathematical audit, 2026-11-14.
% Source: user's DeLellis_Rectifiability.pdf, printed pages cited below.
% This is prose and a specification inventory, not a claim of formal verification.
\chapter{Proof audit: the uniform-measure branch of Preiss' theorem}

\section{Scope, conventions, and conclusion of the audit}
The proposed route through Parts A and B is mathematically appropriate, but the
current graph does not yet specify all the intermediate statements needed for an
implementation. In particular, its moment-coefficient expansion and trace-rigidity
nodes each hide a substantial independent development. The list below separates
these obligations. It is based on Chapters 3 and 6--9 of the supplied book, including
the author's errata at
\url{https://www.math.ias.edu/delellis/sites/math.ias.edu.delellis/files/Errata_libro_0.pdf}.
The errata correct the half-space reduction and the normalized locality theorem,
and explicitly flag the measurability issue in the almost-everywhere recentering
argument. The measurable-hull replacement proposed below is a separate argument
to be proved, not an assertion that the author's suggested projection argument has
already been formalized.

Throughout this audit, $1\leq m\leq n$, measures are positive locally finite Borel
measures on a finite-dimensional real Euclidean space, and
\[
 D^m_{a,r}\mu=r^{-m}(T_{a,r})_\#\mu,
 \qquad T_{a,r}(z)=(z-a)/r,\qquad r>0.
\]
Here powers and scalar multiplication are interpreted in the nonnegative extended
reals when acting on measures. Write $\mathcal U_m$ for measures with
$0\in\operatorname{supp}\nu$ and
$\nu(B(y,r))=r^m$ at every support centre and every $r>0$.
This is the project's unit-ball normalization, not the book's $\omega_m r^m$.
Let $\sigma_V$ be the flat measure on an $m$-plane $V$ normalized by
$\sigma_V(B(0,1))=1$. Equivalently it is a positive scalar multiple of
$\mathcal H^m|V$. Positivity, finiteness and independence of the unit-ball
Hausdorff mass from the choice of plane are explicit normalization lemmas.

Moments divided by their Gaussian mass are unchanged by multiplying the measure
by a positive constant. Consequently the moment identities, trace $m$, and the
rigidity calculations below retain their constants under unit normalization.
The unnormalized Gaussian mass does change: for $\nu\in\mathcal U_m$,
\[
 I_\nu(s)=\int e^{-s|z|^2}\,d\nu(z)
 =\omega_m^{-1}(\pi/s)^{m/2}\quad(s>0),
\]
where $\omega_m$ is the Lebesgue measure of the unit ball in $\mathbb R^m$.
The probability Gaussian on $V$ used in the final trace argument should be
defined by pushing forward the standard $m$-dimensional Gaussian along a linear
isometry, rather than by silently identifying mathlib's Hausdorff normalization
with Lebesgue measure.

\section{Natural-language proof sketch}
At almost every point of positive finite existing density, the existing GMT
theorem makes every tangent uniform. There are tangents: two-sided density
bounds imply local doubling and the upstream compactness theorem then supplies
a nonzero limit. Work at a point where these assertions, and the recentering
property below, all hold.

Part A shows that at least one tangent is flat. Starting with any uniform
tangent, take a tangent at a nearest point to a hole in its support. Its support
lies in a half-space. The corrected barycentre argument then produces either
the same measure or a further tangent supported in some hyperplane. Identify
that hyperplane isometrically with a smaller Euclidean space and repeat.
After finitely many steps a full-dimensional uniform measure is a multiple of
Lebesgue measure. Retain every scaling and isometry witness. A diagonal argument
expresses the result as a limit of rescalings of the original uniform tangent,
with centres in its support. The almost-everywhere recentering property makes
each of these rescalings a tangent of the original measure at the original
point. Closedness of the tangent set then gives a flat tangent.

Part B starts with a flat tangent and excludes a nonflat one. It uses three
facts about unit uniform measures. First, their blow-downs have a unique limit.
Second, this limit is conical, and a conical unit uniform measure sufficiently
close to a plane in an averaged squared-distance sense is flat. Third, a unit
uniform measure with a flat limit at infinity is already flat.

The first fact follows by Gaussian moments. Uniform ball masses identify every
radial integral. A Gaussian translation identity gives a finite Taylor estimate
for moments evaluated at support points. Finite-dimensional graded linear
algebra upgrades these estimates to Taylor coefficients of every moment form.
Hence all Gaussian-weighted polynomial moments of any two blow-down limits
agree. Gaussian tails and a weighted Stone--Weierstrass argument identify the
two measures. Compactness upgrades uniqueness of subsequential limits to
convergence of the whole blow-down family. Scaling uniqueness makes the limit
conical.

For the conical gap, the quadratic Gaussian moment is a positive semidefinite
form of trace $m$, and it equals $|x|^2$ on the support. Its top-$m$ eigenspace
minimizes Gaussian distance energy. If that energy is small, compactness forces
every point of its unit ball to lie near the support. A support point near an
eigenvector corresponding to the $m$th eigenvalue forces that eigenvalue to be
at least one. The trace then forces the form to be an orthogonal projection,
so the measure is supported in an $m$-plane and hence is flat.

For rigidity at infinity, the flat blow-down identifies the leading even
moment coefficients with powers of orthogonal projection onto the limiting
plane. The first support identity gives a normal vector $b$ and
$|Q_Vx|^2=2\langle b,x\rangle$ on the support. Thus the support stays in a bounded
tube about the plane. A second compactness argument gives bounded distance in
the other direction. Fourth-order identities produce a quadratic form
$\widehat b$ on the plane and a quadratic normal-vector map $\omega$.
Gaussian integration gives an upper bound on the second Gaussian moment of
$\widehat b$ in terms of its first moment. The elementary fourth-moment formula
for a Gaussian then yields the sharp trace lower bound. The total trace is also
zero. These facts force $b=0$, so the support is planar and the measure is flat.

Finally, a compactly supported distance energy is continuous both in a measure
and along its scale curve. If a point had both a flat and a nonflat tangent,
the nonflat tangent would have a nonflat blow-down whose energy exceeds the gap
threshold. Between small-energy scales approaching the flat tangent and
large-energy scales approaching that blow-down, choose the last crossing of a
fixed threshold. A subsequence of crossing measures has a uniform limit with
positive energy. The ratio of the nearby flat scales to crossing scales must
diverge; otherwise their limits are related by a finite dilation and both are
flat. Consequently every sufficiently large dilation of the crossing limit
has energy below the threshold. Its limit at infinity is flat by the conical
gap, and rigidity then makes the crossing limit flat, a contradiction.
All tangents are therefore flat almost everywhere. The independently proved
Marstrand--Mattila criterion concludes rectifiability.

\section{Reusable measure and scaling infrastructure}
The following are separate reusable statements, not extra assumptions on Preiss.

\paragraph{S01: compact-test convergence with moving tests.}
If $\mu_i\rightharpoonup\mu$ vaguely, all measures are regular, $K$ is compact,
the continuous tests $f_i,f$ vanish outside $K$, and $f_i\to f$ uniformly,
then $\int f_i\,d\mu_i\to\int f\,d\mu$. First derive a uniform mass bound on
$K$ using a compactly supported cutoff equal to one on $K$.

\paragraph{S02: joint continuity of translated dilations.}
If $\mu_i\rightharpoonup\mu$, $a_i\to a$, $r_i\to r>0$, and
$c_i\to c\in[0,\infty)$, then
$c_i(T_{a_i,r_i})_\#\mu_i\rightharpoonup c(T_{a,r})_\#\mu$.
Assume regularity of the measures and finite real scalars; no bound at infinity
on their total masses is needed. The preimages of a fixed compact test support
lie in one compact set because the parameters converge and $r>0$. The pulled
back tests converge uniformly there. S01 proves the statement. Its special case
$D^m_{a_i,r_i}\mu_i\to D^m_{a,r}\mu$ is used repeatedly.

\paragraph{S03: algebra and covariance.}
Prove
$D^m_{b,s}(D^m_{a,r}\mu)=D^m_{a+rb,rs}\mu$;
positive scalar multiplication commutes with $D$; and pushforward by a linear
isometric embedding intertwines $D$, at centres in its image. Include inverse
dilation and preservation of flatness, regularity and uniform ball constants.

\paragraph{S04: approximation of support points.}
If regular $\mu_i\rightharpoonup\mu$ and $x\in\operatorname{supp}\mu$, then
there exist $x_i\in\operatorname{supp}\mu_i$ with $x_i\to x$, after discarding
finitely many terms if necessary. Use the open-ball portmanteau inequality.
Conversely, if $x_i\to x$, $x_i\in\operatorname{supp}\mu_i$, and the measures
have a common positive lower support-centred ball bound at each radius, then
$x\in\operatorname{supp}\mu$. Bare vague convergence does not give this converse.

\paragraph{S05: stability and compactness of normalized uniform measures.}
A vaguely convergent sequence in $\mathcal U_m$ has a limit in $\mathcal U_m$.
The family is sequentially compact by upstream compact mass bounds.
First approximate support centres using S04; squeeze open and closed radii to
recover the exact ball law, rather than assuming spheres of an arbitrary limit
already have zero mass. The origin's positive ball mass excludes the zero limit.
Also prove the variant where the measures are uniform with the same constant
and their supports approach the origin, without requiring the origin in every
support. This is needed in the bounded-hole argument R05 below.

\paragraph{S06: a countable compact-test metric.}
Construct or reuse a metric on regular positive locally finite Euclidean
measures which detects vague convergence on families with common compact mass
bounds. Include enough cutoff tests to control masses. The necessary outputs are
sequential convergence equivalence, separation, a diagonal selection lemma,
and finite covers of a compact family by metric balls. This can live in a
general vague-convergence module rather than the Preiss theorem file.
For recentering, the finite-cover statement must also cover centered uniform
measures whose ball constants range over a fixed compact interval
$[c,C]\subset(0,\infty)$: density constants at the selected bad centres need
not agree exactly. Unit uniform compactness and continuity of finite scalar
multiplication give this extension.

\paragraph{S07: tangent closure and scaling.}
For a fixed measure and centre, the limit of a vaguely convergent sequence of
dimensional tangents is a dimensional tangent whenever that limit is regular
and nonzero. Prove the analogous arbitrary-scalar statement as a separate
adapter. Constant positive dilations of a dimensional tangent remain dimensional
tangents. At a positive finite density point the dimensional tangent set is
nonempty and compact, because every sufficiently small blow-up has compact
mass bounds and nonzero subsequential limits. A countable-test diagonal lemma
implements closure without assuming a topology on all raw measures.

\paragraph{S08: normalization at a density point.}
Let $q_x=\lim_{r\downarrow0}\mu(B(x,r))/r^m\in(0,\infty)$.
If every arbitrary-scalar tangent is uniform, every dimensional tangent
has ball constant exactly $q_x$ and contains zero in its support.
Use doubling to get the support assertion, the uniform sphere-null property,
and convergence on continuity balls to determine the constant. Dividing the
original measure by $q_x$ therefore puts its dimensional tangents in
$\mathcal U_m$. Existing arbitrary-scalar and dimensional tangent sets are rays
of one another, not literally equal. A pointwise Part B theorem must use this
adapter before applying a universal flatness threshold.

\section{Part A: exact missing statements and measurable-hull repair}
References are Lemma 6.6 and Theorem 6.8, printed pages 57--58 and 63--65;
Corollary 3.9, printed pages 15 and 25--26; and the corrected Lemma 3.8.

\paragraph{A01: support is not full in excess ambient dimension.}
For an $m$-uniform measure on $\mathbb R^k$ with $m<k$, its support is not all
of $\mathbb R^k$. Reuse \texttt{support\_ne\_univ\_of\_lower\_growth}; do
not reconstruct this from the eventual Preiss theorem.

\paragraph{A02: half-space reduction with a retained tangent.}
For uniform $\nu$ with non-full support there exist a support point $y$,
a nonzero vector $e$, and a tangent $\lambda$ at $y$ supported in
$\{z:\langle z,e\rangle\geq0\}$. Preserve the exact normalization when
choosing the tangent. GMT's \texttt{SUniformAux.exists\_halfSpace} proves
this in its body but exports only the existence of the new uniform measure.
The adapter should reuse that proof pattern and retain its witness.

\paragraph{A03: corrected hyperplane reduction with retained lineage.}
For a centred uniform half-space measure $\nu$, there exist $w\ne0$ and a
uniform $\lambda$ supported in $w^\perp$ such that either $\lambda=\nu$
or $\lambda$ is a tangent to $\nu$ at zero, up to the explicitly fixed
positive normalization. The new $w$ need not be $e$.
The proof of GMT's \texttt{SUniformAux.exists\_hyperplane} already has
exactly these two branches. Expose
\texttt{SUniformAux.abs\_inner\_integral\_le},
\texttt{SUniformAux.support\_tangent\_subset\_of\_estimate}, and
\texttt{SUniformAux.support\_subset\_of\_integral\_inner\_eq\_zero}
as upstream prerequisites. The obsolete statement that every tangent lies in
the original half-space boundary must not enter the graph.

\paragraph{A04: hyperplane transport with a map.}
Given an isometric identification $L:\mathbb R^{k-1}\simeq w^\perp$, return
the transported measure together with the equality of its pushforward with
$\lambda$, its ball law, support correspondence and blow-up covariance.
The existential theorem
\texttt{SUniformAux.exists\_isSUniform\_of\_hyperplane}
does not export enough data by itself.

\paragraph{A05: closure of the support-centred scaling orbit.}
Let $\mathcal O_m(\nu)$ be the closure of
$\{D^m_{y,r}\nu:y\in\operatorname{supp}\nu,\ 0<r\leq1\}$ within regular
locally finite measures. For uniform measures all these limits retain their
ball constant. Prove
$\lambda\in\mathcal O_m(\nu)\Rightarrow
 \mathcal O_m(\lambda)\subseteq\mathcal O_m(\nu)$.
To compose a rescaling centred at a support point $z$ of $\lambda$, approximate
$z$ by support points of the convergent approximating measures (S04), then use
S02 and S03. Products of radii remain at most one. A diagonal selection turns
closure membership into the sequence form in the current Lean statement.

\paragraph{A06: a flat measure occurs in that orbit closure.}
Induct on ambient codimension using A01--A05. At dimension $m$, use
\texttt{mattila\_14\_11}, Hausdorff/Lebesgue normalization and isometric
pushforward to identify the terminal uniform measure with a flat one.
There is no need for analytic support varieties, a full Marstrand integer
dimension theorem, or classification of all uniform measures.

\paragraph{A07: measurable hulls preserve disjoint open holes.}
For arbitrary $D$ with $\mu D<\infty$, put $H=\operatorname{toMeasurable}_\mu D$.
Then for every measurable $U$ disjoint from $D$, $\mu(H\cap U)=0$.
This is an immediate use of mathlib's
\texttt{Measure.measure\_toMeasurable\_inter}; an s-finite version is also
available. In particular $\mu(H\setminus\overline D)=0$.
This statement requires no measurability of $D$.

\paragraph{A08: support approximation from an arbitrary positive-mass set.}
Assume $a\in D$, $a$ is a density-one point of its measurable hull $H$,
$\limsup_{r\downarrow0}\mu(B(a,r))/r^m<\infty$, and
$D^m_{a,r_i}\mu\rightharpoonup\nu$ with $r_i\downarrow0$.
Then for every $y\in\operatorname{supp}\nu$,
$\operatorname{dist}(a+r_i y,D)/r_i\to0$.
Indeed, a subsequence of holes $B(a+r_i y,\eta r_i)$ disjoint from $D$ has
zero $H$-mass by A07. It lies in $B(a,Cr_i)$ for a fixed $C$; its remaining
mass is $o(r_i^m)$ by density-one and the upper density bound. Portmanteau
would then make $\nu(B(y,\eta))=0$, a contradiction. Choose approximate
nearest points in $D$ with error $r_i/i$; no closedness of $D$ is needed.

\paragraph{A09: one exceptional set for every recentering.}
Under positive finite existing $m$-density almost everywhere, prove
\[
 \text{for }\mu\text{-a.e. }a,\quad
 \forall\nu\in\operatorname{Tan}_m(\mu,a),\quad
 \forall y\in\operatorname{supp}\nu,\quad\forall r>0,
 \quad D^m_{y,r}\nu\in\operatorname{Tan}_m(\mu,a).
\]
Here $\operatorname{Tan}_m$ uses fixed dimensional normalization.
It is enough to prove recentering at radius one; S03 and S07 give all radii.
For fixed rational tolerances and radius cutoffs, take the arbitrary set of
bad centres, stratify its density constants and intersect with a bounded ball.
If it has positive outer measure, use a finite metric cover of the compact
family of translated uniform tangents to select one positive-outer-measure
class $D$ whose chosen translated tangents are pairwise close.
Take a density point of its measurable hull lying in $D$: the hull's bad
density points are null and cannot cover a set of positive outer measure.
Apply A08 to select nearby bad centres in $D$ and S02 to compare their blow-ups
with the chosen translated tangent. The triangle inequality contradicts the
defining fixed separation of the bad centres. Countably many tolerances give
one null exceptional set. This repairs the book's unjustified assertion that
its existentially defined set $D$ is measurable.

\paragraph{A10: Part A assembly.}
Intersect the a.e. sets for uniformity, density and A09. Choose any tangent,
use A06, recenter every approximating term using A09, and use S07 for the
nonzero flat limit. This is \texttt{ae\_exists\_flat\_tangent}.

\section{Moments and uniqueness at infinity}
References are Lemma 7.2, Definition 7.4, Lemma 7.6 and Proposition 7.7,
printed pages 70--84. All claims below concern positive $s$ and $m>0$.

\paragraph{M01: radial distribution, including integrability.}
For $\nu\in\mathcal U_m$ and $y\in\operatorname{supp}\nu$, identify the
pushforward under $z\mapsto|z-y|$ with that of
$\omega_m^{-1}\mathcal L^m$ under the norm. State nonnegative Borel integrals
as lintegrals, then derive signed Bochner-integral identities only with their
integrability hypotheses. GMT already gives equality of continuous radial
integrals on closed balls. The full pushforward distribution and unbounded
Gaussian consequences are new adapters.

\paragraph{M02: Gaussian growth and uniform tails.}
For every $d\in\mathbb N$ and $s>0$,
$\int |z|^d e^{-s|z|^2}\,d\nu<\infty$ and $0<I_\nu(s)<\infty$.
For fixed $s,d$, the tails outside $B(0,R)$ tend to zero uniformly over
$\mathcal U_m$. State the more reusable version for a family satisfying one
common polynomial bound $\mu(B(0,R))\leq C(1+R)^p$.
Use this to extend vague convergence to polynomial-times-Gaussian tests.
It also justifies differentiation with respect to positive $s$ locally in $s$.

\paragraph{M03: continuous symmetric moment forms.}
Package the current scalar \texttt{gaussianMoment} as a continuous symmetric
$k$-multilinear form, prove the evaluation formula, scalar-invariance,
orthogonal covariance, and the scaling law
\[
 b^{D^m_{0,r}\nu}_{k,s}=r^k b^\nu_{k,s/r^2}.
\]
The definition with $k=0$ is the constant one for uniform measures; Taylor
coefficient assertions below start at $k=1$.

\paragraph{M04: quantitative moment bound and translation identity.}
Prove the norm bound
$|b_{k,s}(u_1,\ldots,u_k)|\leq
C_m 2^k k^{k/2}s^{k/2}\prod_i|u_i|/k!$.
For $x$ in the support, the absolutely convergent series
$\sum_{k\geq0} b_{k,s}(x^k)=e^{s|x|^2}$ follows by translating the
Gaussian. A valid dominating function is
$e^{-s|z|^2+2s|z||x|}$, bounded by a constant times
$e^{-(s/2)|z|^2}$, not the expression with a negative cross term printed in
the source. Taylor truncation gives, for each $q\geq1$,
\[
 \left|\sum_{k=1}^{2q}b_{k,s}(x^k)
 -\sum_{j=1}^{q}\frac{s^j|x|^{2j}}{j!}\right|
 \leq C_{m,q}(s|x|^2)^{q+1/2}.
\]
For $s|x|^2\geq1$ use the finite sum and the norm bound; for small values use
the convergent exponential series with its remainder bound.

\paragraph{M05: analytic weighted projections in a finite grading.}
This is a measure-independent linear algebra theorem. Let the finite-dimensional
real inner-product space $X$ have an orthogonal decomposition
$X=\bigoplus_{j=1}^{N}X_j$, with projections $P_j$, and let $V\leq X$.
For $s>0$ put $A(s)=\sum_j s^jP_j$ and
$F_s=A(s)^{-1}(V^\perp)$. Then $X=V\oplus F_s$; the projection $Q_s$
onto $V$ along $F_s$ extends real-analytically to a neighborhood of zero.
In particular its operator norm stays bounded as $s\downarrow0$.
For the proof, orthogonally split $V^\perp$ according to the filtration
$V^\perp\cap\bigoplus_{j\leq k}X_j$ and define a polynomial family which
is identity on $V$ and equals $P_k+sP_{k-1}+\cdots+s^{k-1}P_1$ on the
$k$th filtration summand. Its value at zero is invertible, by the
least-nonzero-degree argument, and $Q_s=P_V A_s^{-1}$ for that polynomial
family. In the least-degree argument the lower-degree vector belongs to the
sum of all earlier filtration summands, not necessarily only the immediately
preceding one.

\paragraph{M06: support-polynomial space and coefficient extraction.}
For fixed $q$, take the finite sum of symmetric tensor degrees $1$ through
$2q$, with the weighted inner product from Section 7.3, and let $V$ be spanned
by the vectors $(x,x^{\otimes2},\ldots,x^{\otimes2q})$ for support points $x$.
The moment functional $b_s$ vanishes on $F_s$ and M04 determines
$b_s|V$ to order $q$. Choose a finite basis from these spanning support vectors:
its finitely many remainder estimates give an operator-norm
$O(s^{q+1/2})$ bound on $V$. Composing with the bounded $Q_s$ and using M05
gives the Taylor expansion of $b_s$.
Do not implement the source's claim that the unrescaled symmetrized coordinate
monomials are orthonormal: for example
$|e_1\odot e_2|^2=1/2$ for orthogonal unit vectors under its displayed inner
product. Use the inherited tensor inner product or insert multinomial weights.
A finite polynomial-coordinate model is also possible, provided its pairing
realizes $\langle u^{\otimes k},v^{\otimes k}\rangle=\langle u,v\rangle^k$.

\paragraph{M07: coherent Taylor coefficients.}
There exist symmetric continuous forms $b_k^{(j)}$, $k,j\geq1$, such that for
every $q\geq1$,
\[
 b_{k,s}=\sum_{j=1}^{q}\frac{s^j}{j!}b_k^{(j)}+o(s^q),\qquad
 b_k^{(j)}=0\ (k>2j),\qquad
 \sum_{k=1}^{2q}b_k^{(q)}(x^k)=|x|^{2q}\ (x\in\operatorname{supp}\nu).
\]
Uniqueness of finite Taylor coefficients makes choices from different truncation
orders coherent. If $k>2q$, M04 gives $o(s^q)$ directly. The operator-norm
statement is stronger and safer than merely pointwise convergence on the support.

\paragraph{M08: leading scaled moments.}
For odd $k$, $s^{-k/2}b_{k,s}\to0$; for $k=2j$,
$s^{-j}b_{2j,s}\to b_{2j}^{(j)}/j!$.
The factorial is essential; the displayed summary on printed page 74 omits it,
whereas the calculation in Chapter 9 includes it correctly.
Use M03 with $s=r^{-2}$ to deduce the limit of every Gaussian-weighted
polynomial moment of $D^m_{0,r}\nu$ as $r\to\infty$.

\paragraph{M09: Gaussian moment determination.}
Regular measures with a common polynomial growth bound are equal if their
$e^{-\|z\|^2}$-weighted polynomial moments agree.
Weight the two measures by $e^{-\|z\|^2}$. Polynomial growth makes the weighted measures finite. For every vector $w$ and every real $t$, the inequality
\[
t\langle z,w\rangle-\|z\|^2\leq |t|^2\|w\|^2/2-\|z\|^2/2
\]
gives integrability of the exponential of the projection under each weighted measure. The assumed polynomial equalities give equality of all natural moments of every projection. The entire complex moment generating functions have identical derivatives at zero, so analytic uniqueness identifies them. Evaluating at the imaginary unit identifies the characteristic functions of the finite weighted measures. Mathlib's characteristic-function uniqueness theorem identifies these measures, and cancellation of the strictly positive finite Gaussian density identifies the original measures. This avoids expanding a Gaussian at an arbitrary rate without domination.

\paragraph{M10: uniqueness and full convergence at infinity.}
Use S05 to obtain subsequential uniform limits, M02 and M08 to identify their
Gaussian moments, and M09 to identify the limits. The compactness principle
``every sequence escaping to infinity has a convergent subsequence with the
same limit'' gives full convergence. A separate output is that every positive
sequence $r_i\to\infty$ converges to this same measure. Scaling by a fixed
factor preserves the set of limits at infinity, so S02--S03 imply conicality.

\section{The conical gap, without unnecessary classification theorems}
References are Corollary 8.4, Lemmas 8.6--8.7 and Proposition 8.5,
printed pages 85--94. The proof for small energy works for every $m\geq1$;
the stronger unconditional classification in dimensions one and two is not
needed for Preiss and can be omitted. In particular, the two-variable radial
distribution in Lemma 8.6(iii) is not needed by this minimal gap route.

\paragraph{G01: homogeneous conical moments.}
Conicality gives $b_{k,s}=s^{k/2}b_{k,1}$. Comparing with M07 makes odd forms
zero and gives $b_{2j,s}=s^j b_{2j}^{(j)}/j!$.
The support identities give $b_{2j}^{(j)}(x^{2j})=|x|^{2j}$ on the support.

\paragraph{G02: covariance form and its trace.}
The symmetric form $B=b_{2,1}$ is positive semidefinite,
$\operatorname{tr}B=m$, and $B(x,x)=|x|^2$ on the support.
The trace calculation is the radial Gaussian second moment. This statement
does not require spectral choices.

\paragraph{G03: energy comparison for a conical measure.}
For every plane $V$, the distance-squared measure has homogeneous degree
$m+2$, so
$\int e^{-|z|^2}\operatorname{dist}(z,V)^2\,d\nu
=C_m\int_{B(0,1)}\operatorname{dist}(z,V)^2\,d\nu$
for a finite positive constant $C_m$ obtained by radial integration.
With this convention $C_m=\Gamma(m/2+2)$.
State and prove this for nonnegative homogeneous integrands more generally.
It is the missing explicit bridge from the theorem's unit-ball energy to the
Gaussian energy used in the spectral proof.

\paragraph{G04: a spectral minimizing plane.}
For a positive semidefinite symmetric form $B$ on $\mathbb R^n$, choose an
orthonormal eigenbasis with $\alpha_1\geq\cdots\geq\alpha_n\geq0$.
The span $V$ of the first $m$ eigenvectors minimizes
$\operatorname{tr}(B|W^\perp)$ over $m$-planes $W$.
For G02 this trace is a fixed positive multiple of the Gaussian distance
energy. This is a reusable finite-dimensional variational spectral lemma.

\paragraph{G05: small energy implies no holes in the comparison plane.}
For every $\delta>0$ there exists $\eta>0$ such that if $\nu\in\mathcal U_m$
and $V$ is an $m$-plane with Gaussian distance energy at most $\eta$, every
$x\in V\cap\overline B(0,1)$ has distance less than $\delta$ to the support.
Otherwise rotate planes to one fixed plane, take a convergent subsequence of
measures and alleged holes, and use S04--S05. Vanishing energy puts the limit's
support in the plane; the plane-supported uniform classification identifies it
with $\sigma_V$, contradicting the persistent hole. This lemma does not
require conicality or eigenvalue assumptions.

\paragraph{G06: the elementary eigenvalue rigidity step.}
Under G02, assume the top-$m$ eigenspace has a support point within distance
$\delta$ of its $m$th unit eigenvector, where $0<\delta<1/m$.
If $\alpha_m<1$, the trace gives
$\alpha_i-1\leq(m-1)(1-\alpha_m)$ for $i<m$ and
$\alpha_i\leq\alpha_m$ for $i\geq m$.
The support identity then gives
$0\leq(1-\alpha_m)((m-1)^2\delta^2-(1-\delta)^2)<0$.
Thus $\alpha_m\geq1$. Trace $m$ forces the first $m$ eigenvalues to be one
and the rest zero. Hence support is contained in $V$ and $\nu=\sigma_V$.
G03--G06 give the conical gap for all positive $m$ without the flawed
two-dimensional shortcut in the original printed proof.

\section{Rigidity from flatness at infinity: the missing trace chain}
References are Proposition 9.1, Lemmas 9.2, 9.4, 9.9, 9.10 and Proposition
9.5, printed pages 95--109. Assume $\nu\in\mathcal U_m$ has blow-down
$\sigma_V$. Write $P=P_V$, $Q=1-P$, and use the coefficients from M07.

\paragraph{R01: even coefficients determined by the blow-down.}
For every $j\geq1$, $b_{2j}^{(j)}(x^{2j})=|Px|^{2j}$ for every $x$.
Use M02, M03 and the Gaussian moment formula of the flat limit, including the
factorial from M08.

\paragraph{R02: odd top-degree coefficients vanish on the plane.}
For $j\geq1$, the restriction of $b_{2j-1}^{(j)}$ to $V$ is zero.
For a nonzero $v\in V$, use the flat blow-down and S04 to find support points
$x_i$ with $|x_i|\to\infty$ and $x_i/|x_i|\to v/|v|$.
The support polynomial identity, R01 and the fact that all remaining terms have
lower degree show $b_{2j-1}^{(j)}(v^{2j-1})\geq0$.
Repeat with $-v$ and use polarization to obtain the full multilinear restriction.

\paragraph{R03: the normal vector and the cylinder identity.}
Riesz representation and R02 give $b\in V^\perp$ with
$b_1^{(1)}(x)=2\langle b,x\rangle$. The first support identity gives
$|Qx|^2=2\langle b,x\rangle$ on the support, and hence $|Qx|\leq2|b|$.
Also $|Qx-b|=|b|$ there. Keep these identities as named lemmas.

\paragraph{R04: total and normal trace identities.}
The exact radial identity $\operatorname{tr}b_{2,s}=ms$ and M07 imply
$\operatorname{tr}b_2^{(2)}=0$.
R03 and the first-moment expansion give
$\operatorname{tr}(b_2^{(2)}|V^\perp)=8|b|^2
=2\|b_1^{(1)}\|^2$.
These are separate, comparatively short lemmas.

\paragraph{R05: bounded holes in the limiting plane.}
There exists $R_0<\infty$ with
$\operatorname{dist}(v,\operatorname{supp}\nu)<R_0$ for every $v\in V$.
If not, choose $v_i$ with distances $r_i\to\infty$, a nearest support point,
and its projection $z_i\in V$. The measures $D^m_{z_i,r_i}\nu$ have
supports approaching the origin and lie in shrinking tubes about $V$ by R03.
S05 gives a unit uniform limit supported in $V$, hence $\sigma_V$.
The original holes give a unit hole centred at a bounded point of $V$, a
contradiction. The centres $z_i$ need not lie in the support, which is why the
near-support variant of S05 is essential.

\paragraph{R06: quadratic normal map and adjusted quadratic form.}
Define the symmetric bilinear normal-valued map $\omega$ by
\[
 \langle\omega(u,u),w\rangle
 =3b_3^{(2)}(u,u,w)-4|u|^2\langle b,w\rangle\quad(u\in V),
\]
and put $\widehat b(u,u)=b_2^{(2)}(u,u)+\langle\omega(u,u),b\rangle$.
Use polarization for mixed arguments and R02 to prove $\omega$ takes values
in $V^\perp$. These definitions need actual typed bilinear maps, not just a
scalar diagonal moment notation.

\paragraph{R07: the Gaussian limit and its vanishing trace contraction.}
For $v,w\in\mathbb R^n$, prove
\[
 \lim_{s\downarrow0}\frac{8s}{I(s)}\int e^{-s|z|^2}
 \langle z,v\rangle^2\langle z-b,w\rangle\,d\nu
 =3b_3^{(2)}(v,v,w)-4b_2^{(1)}(v,v)\langle b,w\rangle.
\]
For $v\in V,w\in V^\perp$ this is $\langle\omega(v,v),w\rangle$.
Further prove
\[
 \lim_{s\downarrow0}\frac{s}{I(s)}\int e^{-s|z|^2}
 |Pz|^2\langle z-b,w\rangle\,d\nu=0.
\]
To justify the latter, the analogous $|z|^2$ expression has a limit by the
first formula and a finite orthonormal-basis sum. Differentiating the Gaussian
first moment and applying the appropriate one-sided l'Hopital argument
identifies this limit as zero; the $|Qz|^2$ term tends to zero by R03 and M02.
Record differentiability and domination hypotheses before using l'Hopital.

\paragraph{R08: Gaussian trace formula.}
Let $\gamma$ be the standard covariance-one Gaussian probability measure on
$V$. Prove $\int\langle\omega(v,v),b\rangle\,d\gamma(v)=0$ using R07,
Fubini and domination by a constant times $|v|^2$.
Then
$\operatorname{tr}(b_2^{(2)}|V)=\int\widehat b(v,v)\,d\gamma(v)$.
The same argument works against each fixed normal vector, not only $b$.

\paragraph{R09: vanishings and the fourth-degree support identity.}
For $v\in V,w\in V^\perp$,
$b_3^{(2)}(v,v,v)=b_3^{(2)}(v,w,w)=b_3^{(2)}(w,w,w)=0$.
The first is R02; the others follow by R03's bounded normal coordinates and
Gaussian estimates in the defining limits. Therefore, on the support,
\[
 b_1^{(2)}(z)+b_2^{(2)}(z,z)+3b_3^{(2)}(Pz,Pz,Qz)
 =|Qz|^4+2|Pz|^2|Qz|^2.
\]
Use the order-two support identity and R01; do not drop mixed cubic terms
without first establishing these vanishings.

\paragraph{R10: a pointwise inequality for the adjusted form.}
For all $v\in V$,
$\widehat b(v,v)^2\leq |b|^2|\omega(v,v)|^2$.
Choose support points $t_i v+v_i+w_i$ within bounded distance of $t_i v$ using
R05. Pass to a convergent subsequence of bounded normal components $w_i$.
Divide R09 by $t_i^2$ and use R03 to obtain
$b_2^{(2)}(v,v)+\langle\omega(v,v),w\rangle=0$ and
$|w|^2=2\langle b,w\rangle$.
Thus $|w-b|=|b|$ and
$\widehat b(v,v)+\langle\omega(v,v),w-b\rangle=0$; Cauchy--Schwarz proves
the inequality.

\paragraph{R11: lower-order error vanishes after Gaussian rescaling.}
R09 and R03 imply that
$\widehat b(Pz,Pz)+\langle\omega(Pz,Pz),Qz-b\rangle$
has absolute value at most $C(1+|z|)$ on the support.
Hence its integral against $sI(s)^{-1}e^{-s|z|^2}\nu$ tends to zero.
This is equation (9.34); it is a necessary bridge to the global trace bound.

\paragraph{R12: fourth Gaussian moments and a reproducing identity.}
For $y,z\in V$,
$\int\langle y,v\rangle^2\langle z,v\rangle^2\,d\gamma
=|y|^2|z|^2+2\langle y,z\rangle^2$.
The implemented Gaussian foundation first differentiates the real Gaussian
moment generating function $\exp(\sigma^2t^2/2)$ at zero. Leibniz's rule and
$f'(t)=\sigma^2tf(t)$ give every even moment, including zero variance.
Linear pushforward identifies each projection of the standard Gaussian with
this real Gaussian. Polarizing the fourth power gives the displayed mixed
moment; a second polarization gives the four-projection Wick formula.
All finite products of projections are integrable by Fernique's norm-moment
bounds and the inequality
$\prod_i|\langle u_i,v\rangle|\leq(\prod_i|u_i|)|v|^{\#i}$.
Combining this with R07 yields, for $w\in V^\perp$,
\[
 \int\langle y,v\rangle^2\langle\omega(v,v),w\rangle\,d\gamma(v)
 =2\langle\omega(y,y),w\rangle.
\]
Justify limits by Gaussian quartic domination and Fubini.

\paragraph{R13: the integral quadratic inequality.}
Integrate R10, express $\int|\omega(v,v)|^2\,d\gamma$ through R07 and R12,
use R11, and pass to the flat blow-down using M02. The result is
\[
 \int\widehat b(v,v)^2\,d\gamma(v)
 \leq-8|b|^2\int\widehat b(v,v)\,d\gamma(v).
\]
The inequality must remain an inequality throughout. The equality sign at
the start of the second long display on printed page 108 loses the earlier
inequality; its final statement (9.41) has the correct direction.

\paragraph{R14: abstract Gaussian quadratic-form bound.}
For every symmetric bilinear form $A$ on an $m$-dimensional Euclidean space,
\[
 \int A(v,v)^2\,d\gamma=(\operatorname{tr}A)^2+
 2\|A\|_{\mathrm{HS}}^2
 \geq(1+2/m)\left(\int A(v,v)\,d\gamma\right)^2.
\]
This is a reusable Gaussian theorem independent of uniform measures.
The implemented proof expands the quadratic form in an arbitrary orthonormal
basis, integrates its square using the four-projection Wick formula, and
uses Parseval to identify the coefficient-square sum. Cauchy--Schwarz
on the diagonal coefficients gives the dimension-dependent lower bound.
Combining it with R13 gives
$\int\widehat b(v,v)\,d\gamma\geq-8m|b|^2/(m+2)$.
R08 and R04 therefore give
$\operatorname{tr}b_2^{(2)}\geq16|b|^2/(m+2)$.
Since that trace is zero, $b=0$; R03 makes the support planar, and plane-supported
uniform classification finishes rigidity. This is the actual content hidden
inside the current single trace-rigidity node.

\section{First-crossing assembly and its exact quantifiers}
References are Lemma 6.20 and Section 6.4, printed pages 66--69.

\paragraph{C01: compact projection model for planes.}
Use self-adjoint idempotent operators of trace $m$ to model $m$-planes.
This is a nonempty compact set when $m\leq n$.
Its distance functions converge uniformly on bounded sets when the projections
converge. No lattice topology on the submodule type is intended.

\paragraph{C02: a nonnegative cutoff and energy.}
Choose $\phi\in C_c(\mathbb R^n)$ with $0\leq\phi\leq1$,
$\phi=1$ on $\overline B(0,1)$ and support in $B(0,2)$.
Define $F(\mu)=\min_V\int\phi(z)\operatorname{dist}(z,V)^2\,d\mu$.
The minimum is attained. S01 and C01 imply sequential continuity for vague
convergence; S02 implies continuity of $r\mapsto F(D^m_{x,r}\mu)$ for $r>0$.
Flat measures have $F=0$. For a unit uniform measure, $F\leq\varepsilon$
implies that some plane has unit-ball energy at most $\varepsilon$.

\paragraph{C03: a blow-down of a tangent is still a tangent.}
If a dimensional tangent $\zeta$ has a nonzero regular limit at infinity $\chi$,
then $\chi$ is a dimensional tangent of the original measure at the same point.
Use S07's positive dilation invariance and closure. This statement is pointwise;
it needs no almost-everywhere recentering theorem.

\paragraph{C04: first-crossing radii with an interval bound.}
After S08 normalization, suppose $\alpha$ is a flat dimensional tangent and
$\zeta$ is a nonflat dimensional tangent, all dimensional tangents being unit
uniform. Let $\chi$ be the blow-down of $\zeta$. Rigidity and the gap give
$F(\chi)>\varepsilon$. Choose $r_i\to0$ with blow-ups tending to $\alpha$
and $s_i<r_i$ with blow-ups tending to $\chi$; subsequence selection makes
the ordering possible. For sufficiently large $i$, the continuous scale
function satisfies $F(D_{x,r_i}^m\mu)<\varepsilon$ and
$F(D_{x,s_i}^m\mu)>\varepsilon$.
Choose $\sigma_i\in[s_i,r_i]$ with energy exactly $\varepsilon$ and energy
at most $\varepsilon$ on the entire interval $[\sigma_i,r_i]$.
This is a greatest-crossing/last-crossing assertion, not merely the intermediate
value theorem. Extract a nonzero unit uniform limit $\xi$ at the $\sigma_i$.

\paragraph{C05: scale ratios diverge.}
Prove $r_i/\sigma_i\to\infty$. If the ratios have a bounded subsequence,
extract a further limit $R\in[1,\infty)$. S02--S03 give
$D^m_{0,R}\xi=\alpha$. Invert the dilation to conclude $\xi$ flat,
contradicting $F(\xi)=\varepsilon>0$.

\paragraph{C06: the crossing limit is flat at infinity.}
For every fixed $R\geq1$, eventually $\sigma_i R\in[\sigma_i,r_i]$.
S02--S03 and C04 give $F(D^m_{0,R}\xi)\leq\varepsilon$.
Let $R\to\infty$. M10 and continuity give $F(\psi)\leq\varepsilon$ for
the conical blow-down $\psi$ of $\xi$. Apply the conical gap, then rigidity,
to conclude $\xi$ flat, the final contradiction.
Keep the centre $x$ fixed throughout; the printed scale function sometimes
uses zero where $x$ is intended. Every use of $F$ on a dilated measure must
retain the factor $R^{-m}$.

\section{Suggested dependency and module boundaries}
The key paths are
\[
 \text{S01--S08 + A01--A09}\longrightarrow\text{Part A},\qquad
 \text{M01--M09}\longrightarrow\text{M10},
\]
\[
 \text{M07 + M10 + G01--G06}\longrightarrow\text{conical gap},\qquad
 \text{M07 + M10 + R01--R14}\longrightarrow\text{rigidity},
\]
\[
 \text{S07--S08 + M10 + gap + rigidity + C01--C06}
 \longrightarrow\text{Part B}.
\]
Part A and Part B then feed the independent Marstrand--Mattila criterion.

S01--S06 belong in general vague-convergence and measure-pushforward modules;
S07--S08, A05 and A07--A09 belong in tangent-measure modules. The corrected
uniform-reduction adapters A01--A04 should stay close to the existing GMT
Marstrand development. M05 belongs in finite-dimensional graded linear algebra;
M02 and M09 can be general Gaussian-weighted-measure theorems; M03--M08 and
G01--G02 are uniform-moment modules. C01 is Euclidean Grassmannian geometry,
G04 finite-dimensional spectral optimization, G05 uniform compactness,
R14 Gaussian quadratic-form integration. Only the synthesis of those results
belongs in files called Preiss Part A, Part B or Main.

The current scalar moment definition is an appropriate initial evaluation API,
but cannot by itself state tensor Taylor expansions, restrictions, traces,
polarization, $\omega$, or $\widehat b$. Those objects must be designed before
the blueprint can be called a complete set of Lean interfaces. The generic
graded projection theorem is a preferable reusable abstraction to formalizing
all of the book's symmetric-tensor notation literally.

\paragraph{A concrete Lean-friendly alternative to symmetric tensor quotients.}
For the public API, use
\texttt{ContinuousMultilinearMap R (fun \_ : Fin k => E) R}
with a theorem asserting invariance under permutations of the input tuple.
Coefficient families have type
\texttt{(k j : Nat) -> ContinuousMultilinearMap R (fun \_ : Fin k => E) R}.
State their remainder with \texttt{Asymptotics.IsLittleO} along the positive
neighborhood filter at zero, using the operator norm of continuous multilinear
maps. State the diagonal support identities separately. A bilinear form can be
curried into a continuous linear map to the dual, then represented by a
self-adjoint endomorphism via Riesz. Define its trace using the existing
\texttt{LinearMap.trace}; independence of orthonormal basis and the diagonal
sum formula are reusable lemmas, rather than definitions tied to one coordinate
basis. For $\omega$, use a continuous bilinear map
$V\to V\to V^\perp$; for $\widehat b$, a continuous symmetric bilinear form
on the subtype $V$.

The proof of M06 can use only finite Euclidean coordinate spaces, without
constructing a symmetric-tensor quotient. Index a coordinate block of degree
$k$ by all functions $\iota:\operatorname{Fin}k\to\operatorname{Fin}n$, and
define
\[
 \Phi_k(x)_\iota=\sqrt{2^k/k!}\prod_{j=1}^k x_{\iota(j)}.
\]
The finite-sum/product identity gives
$\langle\Phi_k(x),\Phi_k(y)\rangle=(2^k/k!)\langle x,y\rangle^k$.
Take the Euclidean space indexed by the finite disjoint union of blocks
$1\leq k\leq2q$ and put $\Phi=(\Phi_1,\ldots,\Phi_{2q})$.
The coordinate-block projections are an orthogonal grading. Its moment
functional is
\[
 u\longmapsto I(s)^{-1}\int
 \left\langle\left(\sum_{k=1}^{2q}s^kP_k\right)u,\Phi(z)\right\rangle
 e^{-s|z|^2}\,d\nu(z).
\]
Evaluated at $\Phi(x)$, it is exactly $\sum_{k=1}^{2q}b_{k,s}(x^k)$.
For the subspace spanned by $\Phi(\operatorname{supp}\nu)$ this functional
vanishes on the weighted complement in M05. The target polynomial
$\sum_{j=1}^q s^j|x|^{2j}/j!$ is a linear functional of $\Phi(x)$, since each
homogeneous norm polynomial has an explicit coordinate expansion. M05 now
applies verbatim. Extract coordinate coefficients and reconstruct the
multilinear forms by finite sums; permutation symmetry follows from the
symmetry of the original Gaussian integral and uniqueness of Taylor
coefficients. The redundant nonsymmetric coordinates cause no problem: the
support-generated subspace and its orthogonal complement handle them.
This design uses \texttt{EuclideanSpace}, finite sums, submodules and existing
continuous linear/multilinear maps; its only new abstract theorem is M05.

The current graph correctly distinguishes unit uniform measures from arbitrary
positive scalar uniform measures, excludes zero from tangents, and keeps the
moving-centre pseudotangent theorem separate from ordinary tangents. Preserve
these distinctions. Merge the two consecutive proof environments currently
attached to its radial-integral node: the graph checker associates only the
first one with the statement.
